\chapter{Conclusão}\label{cap:conclusao}

Esta dissertação avaliou a transformação de uma reconstrução neural com decomposição de instâncias em uma representação explícita e estruturada para simulação robótica. A cadeia manteve o DM-NeRF como método neural e reuniu extração geométrica, individualização dos elementos, regularização de uma superfície de apoio e integração com o Gazebo Classic.

Na cena \textit{study}, a avaliação oficial em 100 vistas obteve PSNR de 41,1100 dB, SSIM de 0,9875 e LPIPS de 0,0260. A decomposição apresentou AP50 de 0,9956 e AP95 de 0,7192. A malha extraída atingiu F-score de 0,0480 a 1 cm, 0,1960 a 2 cm e 0,7750 a 5 cm, evidenciando a diferença entre qualidade de síntese e correspondência geométrica. No mesmo suporte, o COLMAP obteve 0,0285, 0,1288 e 0,3721 nesses limiares.

A separação estrutural gerou 13 objetos, preservou 95,844\% das faces homogêneas e registrou 29.545 faces de fronteira. A regra geométrica selecionou a instância 10 como superfície de apoio, e sua regularização reduziu o P95 do residual ao plano de 6,80 cm para erro numérico. Essas etapas estabeleceram uma correspondência explícita entre a topologia reconstruída, os identificadores de instância e os modelos exportados.

No Gazebo Classic, as três condições foram carregadas, uma instância foi manipulada de forma independente e a superfície regularizada sustentou o corpo dinâmico no ponto ensaiado. O protocolo distinguiu, assim, carregamento, independência estrutural e contato, em vez de tratar a simples abertura da malha no simulador como evidência suficiente.

Os resultados respondem afirmativamente à questão de pesquisa para a cena \textit{study} e o protocolo adotado: o pós-processamento transformou a reconstrução DM-NeRF em geometria explícita, objetos independentes e uma cena funcional no simulador. A contribuição do trabalho está na caracterização e avaliação dessa passagem, com evidências separadas para aparência, geometria, estrutura e funcionamento.
